\documentclass{article}
\usepackage[T1]{fontenc}
\usepackage{lmodern}
\usepackage{graphicx}
\usepackage{amsmath,amssymb}
\usepackage[a4paper,margin=2cm]{geometry}
\usepackage[absolute,overlay]{textpos}
\setlength{\TPHorizModule}{1mm}
\setlength{\TPVertModule}{1mm}
\usepackage[indonesian]{babel}
\usepackage{hyperref}
\setlength{\emergencystretch}{2em}
% unit: Halaman 10

\begin{document}

% === Halaman 10 (llm) ===

\begin{itemize}
    \item \textbf{Contoh 1:}

    Plainteks: \texttt{onetimepad}\
    Kunci: \texttt{tbfrgfarfm}

    Misalkan $A = 0, B = 1, ..., Z = 25$.

    cipherteks: \texttt{HOJKOREGHP}

    yang dalam hal ini diperoleh sebagai berikut:
    \[
    \begin{aligned}
    ('o' + 'T') \pmod{26} &= H \\
    ('n' + 'B') \pmod{26} &= O \\
    ('e' + 'F') \pmod{26} &= J, \text{ dst}
    \end{aligned}
    \]
\end{itemize}

\end{document}
